\section{Background and Related Work}
\label{sec:background}

Four communities have addressed pieces of the problem
\texttt{MaybeUncertain<T>} solves. None has combined them into a value-level
type with statistically-gated lifting. We survey each in turn and state what
each does not provide.

\subsection{The \texttt{Uncertain<T>} lineage}
\label{sec:bg:uncertain}

Bornholt et al.~\cite{bornholt2014uncertain,bornholt2013thesis}
introduced \texttt{Uncertain<T>} at ASPLOS 2014. The type encapsulates a
random variable of a numeric type $T$ and overloads arithmetic, comparison,
and logical operators to construct a Bayesian network of computations.
Conditionals on \texttt{Uncertain<bool>} are discharged by Wald's sequential
probability ratio test~\cite{wald1945sprt,wald1948optimum}, which
draws samples until the evidence supports rejecting either the null or the
alternate hypothesis at the chosen confidence level. McKinley's POPL 2016
keynote~\cite{mckinley2016popl} retrospectively framed
\texttt{Uncertain<T>} as a programming-language abstraction for uncertain data
and invited extensions in other languages.

The companion line at the same group includes probabilistic
assertions~\cite{sampson2014passert}, which verify quantitative properties of
randomised programs via distribution extraction, and
EnerJ~\cite{sampson2011enerj} and Rely~\cite{carbin2013rely}, which apply
type-level reasoning to approximate computing. Stanley-Marbell et al.'s recent
survey~\cite{stanley2025approxsurvey} positions \texttt{Uncertain<T>} as the
leading programming-language abstraction in this space and notes no presence
extension among its descendants.

Community ports include a Swift implementation by Thompson~\cite{mattt2019swiftuncertain}
and an independent Rust port (\texttt{uncertain-rs})~\cite{uncertainrs}. Neither
introduces a presence factor. \texttt{MaybeUncertain<T>} is implemented in the
\texttt{deep\_causality\_uncertain} crate of the DeepCausality
project~\cite{hansen_deepcausality} and is, to our knowledge, the first
value-level type that factors probabilistic presence from value uncertainty.

\paragraph{What this lineage does not provide.} The presence of the value is
assumed. Missing data is silently absorbed into a degenerate distribution at
the value itself, conflating measurement noise with the failure to measure.

\subsection{Probabilistic programming languages}
\label{sec:bg:ppl}

Probabilistic programming languages (PPLs) form a large
landscape~\cite{tolpin2016anglican,carpenter2017stan,bingham2019pyro,ge2025turing,cusumano2019gen,narayanan2016hakaru,holtzen2020dice,gehr2016psi,gehr2020lambdapsi,milch2005blog,kimmig2011problog,pfeffer2009figaro}.
The relevant question for this paper is narrow: do any of them model the
optional-ness of a value as a type-level concept distinct from the value's
distribution? The closest mechanical analogue is Pyro's
\texttt{MaskedDistribution} and the \texttt{pyro.mask} handler~\cite{bingham2019pyro},
which let users zero out the log-density contribution of observations during
inference. This is a runtime decorator on a distribution object, not a value
type. The masked observation is still numerically present in the
data structure; the mask only affects how the inference engine accounts for it.

Stan models missing data by promoting absent values to additional parameters
of the model~\cite{carpenter2017stan}. The user writes the missingness model
by hand. Turing.jl~\cite{ge2025turing} forwards Julia's \texttt{Missing} union
through the inference pipeline and treats missing observations as latent
variables. In each case the abstraction lives at the model layer; the
programmer writes presence into the model, rather than receiving a value-level
type that already encodes it.

\paragraph{Probabilistic program verification logics.} Recent work has
extended type systems and Hoare logics with relational and probabilistic
reasoning. \emph{Safe Couplings}~\cite{vasilenko2022safecouplings} lifts
refinement types in Liquid Haskell to relational probabilistic properties via
couplings, with case studies on Temporal Difference learning and Stochastic
Gradient Descent stability. eRHL~\cite{avanzini2025erhl} is a quantitative
relational Hoare logic with soundness and completeness for almost-surely
terminating programs and characterises statistical distance and differential
privacy. Both occupy the verification-logic layer above the value-abstraction
layer where \texttt{MaybeUncertain<T>} lives; neither targets the
presence-vs-value factoring.

\paragraph{What PPLs do not provide.} A first-class value-level type that
encapsulates presence-uncertainty and composes through ordinary arithmetic
operators is not present in any PPL we surveyed. The factoring exists at the
inference layer and at the modelling layer; it does not exist at the type
layer.

\subsection{Missing-data statistics}
\label{sec:bg:missingdata}

Rubin's 1976 taxonomy~\cite{rubin1976missing} distinguishes three
missingness mechanisms. Data is \emph{missing completely at random} (MCAR) if
the probability of missingness is independent of both observed and unobserved
values. Data is \emph{missing at random} (MAR) if missingness depends only on
observed values. Data is \emph{missing not at random} (MNAR) if missingness
depends on the unobserved value itself. The standard
text~\cite{littlerubin2019} develops estimators for each regime and the
boundary cases~\cite{mealli2015mar,seaman2013mar} clarify when the
distinctions matter for inference. Recent machine-learning work treats MNAR
directly~\cite{ipsen2021notmiwae}.

\texttt{MaybeUncertain<T>}'s independence assumption between $\mathit{is\_present}$
and $\mathit{value}$ corresponds to MAR semantics. We state this assumption
explicitly in~\Cref{sec:type} and revisit its limitations in~\Cref{sec:discussion}.

\paragraph{What the statistics literature does not provide.} A programming
abstraction. The literature gives the semantics; \texttt{MaybeUncertain<T>}
provides the type that implements it.

\subsection{Adjacent constructions: probabilistic databases and sensor fusion}
\label{sec:bg:adjacent}

Two communities have already factored presence from value in their own
settings.

\emph{Probabilistic databases} model tuple-existence probability separately
from attribute-value uncertainty. ULDBs~\cite{benjelloun2006uldb} attach
existence probabilities to tuples and propagate lineage through queries.
MayBMS~\cite{antova2007maybms} uses world-set decompositions that formally
separate ``the tuple exists'' from ``the attribute has this value.''
Orion~\cite{cheng2008orion} extends relational systems with both
tuple-level and attribute-level uncertainty. The standard
text~\cite{suciu2011pdb} synthesises the field. The factoring is structurally
identical to what \texttt{MaybeUncertain<T>} performs at the value level, but
occurs at the relational layer; it does not give the application programmer a
value-level type usable anywhere \texttt{Option<T>} would go.

\emph{Sensor fusion and intermittent observations} factor presence from
value inside numerical filters. Sinopoli et al.~\cite{sinopoli2004kalman}
analyse Kalman filtering when measurements arrive according to a Bernoulli
process; subsequent work~\cite{liu2004kalman,censi2011kalmansemi} refines the
convergence theory, and out-of-sequence-measurement
methods~\cite{orton2005oosm,barshalom2004oosm} treat asynchronous arrivals.
The factoring is mathematically the same Bernoulli-times-Gaussian construction
\texttt{MaybeUncertain<T>} implements, but lives inside a specific numerical
filter rather than as a general-purpose type.

\paragraph{What these communities do not provide.} A value-level type usable
in arbitrary application code. Probabilistic databases solve the conceptual
problem at the relation level; sensor-fusion methods solve it inside a filter.
Neither exports the factoring to ordinary host-language values.

\subsection{Effect systems and the probability monad}
\label{sec:bg:effects}

The categorical foundations are well established. Giry's measure
monad~\cite{giry1981categorical} and Kozen's probabilistic dynamic
logic~\cite{kozen1985probpdl} provide the semantic ground. Algebraic effects
and handlers~\cite{plotkin2009handlers,bauer2015eff,leijen2014koka} give a
compositional framework in which probability and partiality can be combined,
and graded or parametric effect monads~\cite{katsumata2014parametric}
formalise the stack structure that a \texttt{Maybe}-over-probability
construction inhabits. Statistical model checking~\cite{younes2006sprtmc} uses
SPRT inside verification tools, demonstrating SPRT's relevance beyond
Bornholt's setting.

\texttt{MaybeUncertain<T>} sits adjacent to this literature without
instantiating any specific construction from it. The implementation provides
no monadic combinators; no \texttt{bind} or \texttt{flat\_map} operation is
exposed, and the type does not implement the \texttt{Monad},
\texttt{Functor}, or \texttt{Applicative} traits available in the sibling
\texttt{deep\_causality\_haft} crate. Formalising the type as a (graded)
monad with stated and proved laws is left for future work. Probabilistic
affine forms~\cite{bouissou2016probaffine} and classical interval
analysis~\cite{moore1966interval} offer guaranteed enclosures rather than
sampled distributions and are complementary rather than competing.

\paragraph{What this layer does not provide.} A specific, implementable type
for a working programming language. The effect-system literature gives the
algebra; \texttt{MaybeUncertain<T>} instantiates it as a Rust type with
operators, an SPRT-gated bridge to \texttt{Uncertain<T>}, and measured
overhead.

\subsection{Summary}

\texttt{MaybeUncertain<T>} occupies a small gap in a well-populated landscape.
The \texttt{Uncertain<T>} lineage provides the value-uncertainty type and the
SPRT machinery but not the presence axis. PPLs handle missing data at the
modelling or inference layer but not at the type layer. Statistics gives the
semantics but not the abstraction. Probabilistic databases and sensor fusion
factor presence from value but at the wrong abstraction layer for a
general-purpose program. The effect-system literature gives the algebra but
not a concrete usable type. \texttt{MaybeUncertain<T>} fills the gap with a
small, principled extension of \texttt{Uncertain<T>}: one Bernoulli node per
leaf, AND-propagation on operators, one gating operator built on the existing
SPRT.
